\subsection{Principal findings}
Logistic regression achieved the strongest pooled LOPO discrimination, while random forest achieved higher accuracy, sensitivity, and F1. Stacking offered the highest specificity but did not improve AUC over logistic regression. These trade-offs indicate that added complexity does not confer a uniform advantage and that probability ranking and classification at threshold 0.5 are distinct objectives. The observed differences do not establish statistically significant superiority. Logistic regression remains a useful interpretable baseline, while selection of an operating model depends on the relative consequences of missed cases and false positives.

The findings accord with literature showing that optimisation and ensemble complexity do not consistently improve performance on Gallstone-1. Published scores cannot be compared directly with these LOPO estimates because training sizes, feature-selection procedures, thresholds, and assessment samples differ. Likewise, errors recurring across model families suggest shared limitations in the available measurements or fitted relationships, but retrospective inspection cannot determine a record's clinical circumstances or establish incorrect labels.

\subsection{Limitations}
The analysis relies on a small cohort from one clinical setting. Its near-balanced case distribution does not establish predictive values or assessment workload in populations with different prevalence. Availability of detailed bioimpedance measurements may further constrain transferability. Source-row identifiers are not verified person identifiers; the LOPO interpretation assumes one record per person, and repeated-person data would require grouped exclusion.

Fold-specific preprocessing prevents direct use of excluded records during fitting, but the analysis evaluates fixed pipelines rather than the entire model-development process. The dataset and modelling choices had been explored previously; consequently, LOPO is an internal comparison rather than an untouched confirmatory assessment. Retrospective selection of the strongest family can introduce optimism, and the heavily overlapping training folds are not independent experiments. The single fixed-seed run does not quantify random-state variability or furnish confidence intervals for model differences.

The study did not establish calibrated probabilities, optimise a clinical operating threshold, or demonstrate external generalisation or clinical benefit. Feature selection may also vary across folds because correlated measurements can substitute for one another. Aggregate metrics and recurring error examples therefore require cautious interpretation and do not establish causal relationships.

\subsection{Implications for further research}
Further work should place feature, parameter, calibration, threshold, and model-family decisions inside a nested development procedure. A sensitivity-oriented threshold should be assessed jointly with specificity and assessment workload. Grouped feature ablation and prespecified subgroup evaluation could investigate measurement burden and uneven errors. An untouched external cohort is required to assess transfer beyond this dataset. Prospective disease-onset prediction would require a different study design and outcome definition.
